\section{Whole-tensor parametrizations}\label{atlas:section}

We give the whole tensors in each of the nine types.  In this section
\((l,m,t,z)\) is an ordering of the four sites, and
\[
 L=W_{lm},\quad P=W_{lt},\quad Q=W_{lz},\quad
 R=W_{mt},\quad S=W_{mz},\quad K=W_{tz}.
\]
Juxtaposition denotes tensor product, restored without a sign to the indicated
site order.  The equation to be solved is
\begin{equation}\label{atlas:equation}
 LK+PS+QR=\tau c_lc_mc_tc_z,\qquad \tau\ne0.
\end{equation}
Write \(V_i=\C c_i\oplus H_i\), using the fixed complements to identify
quotients with \(H_i\).  A bar on a vector is its quotient image; on an edge
tensor it means projection at both endpoints.  In particular,
\begin{equation}\label{atlas:zero-kernel}
 \overline{W}_{ij}=0
 \quad\Longleftrightarrow\quad
 W_{ij}\in c_i\otimes V_j+V_i\otimes c_j.
\end{equation}
A tensor described as free with prescribed projected rank may have arbitrary
entries involving a \(c\) factor, including entries making its whole matrix
singular.  All scalar and vector parameters below are complex.

\subsection{Empty projected support}

Let \(E=\C^6\), with coordinates indexed by the six unordered pairs of sites,
and put
\[
 q(x)=x_{12}x_{34}+x_{13}x_{24}+x_{14}x_{23},\qquad
 b(x,y)=q(x+y)-q(x)-q(y).
\]
Let \(S_i\subset E\) be the coordinate subspace on the three edges incident
to \(i\).

\begin{proposition}\label{atlas:empty}
The solutions with all six double projections zero are parametrized as follows.
Choose \(a\in E\) with \(q(a)=\tau\), a totally isotropic subspace
\(U\subset a^\perp\) for \(b\), and linear maps
\[
 B_i:H_i^*\longrightarrow U\cap S_i\qquad(1\le i\le4).
\]
Necessarily \(\dim U\le2\).  Define \(u_{ij}\in H_i\) by
\((B_i(x))_{ij}=x(u_{ij})\), and set
\begin{equation}\label{atlas:empty-form}
 W_{ij}=a_{ij}c_ic_j+u_{ij}c_j+c_iu_{ji}.
\end{equation}
These formulas give every solution of this type.  The data \(a,B_i\) are
unique; the choice of \(U\) becomes unique if
\(U=\sum_i\im B_i\) is required.
\end{proposition}

\begin{proof}
By \eqref{atlas:zero-kernel}, every such tuple has the unique form
\eqref{atlas:empty-form}.  Evaluate the tensor at covectors taking \(c_i\)
to \(1\), and denote their restrictions to \(H_i\) by \(x_i\).
The six edge evaluations form the vector
\(a+\sum_i B_i(x_i)\).  Thus the target identity is equivalent to
\[
 q\left(a+\sum_iB_i(x_i)\right)=\tau
 \quad\text{for all }x_i\in H_i^*.
\]
This equivalence loses no coefficients: the expression is multiaffine in
the four groups of variables, and homogenizing each group recovers the
original four-tensor.  Comparison of constant, linear and quadratic terms,
followed by polarization, gives precisely
\[
 q(a)=\tau,\qquad b(a,U)=0,\qquad b(U,U)=0,
 \qquad U=\sum_i\im B_i.
\]
The quadratic terms give \(b(\im B_i,\im B_j)=0\) for \(i\ne j\); for
\(i=j\) this holds because each \(S_i\) is totally isotropic, no two
edges at \(i\) being complementary.
Conversely these conditions make the displayed polynomial constant.
The form \(b\) is nondegenerate, since its matrix has three invertible
complementary-edge blocks.  As \(b(a,a)=2\tau\ne0\), the space \(a^\perp\)
is nondegenerate of dimension five.  For an isotropic \(U\) in that space,
\(U\subset U^\perp\), so \(2\dim U\le5\).
\end{proof}

\subsection{A single projected edge}

We first record the crossing decomposition needed at a zero opposite edge.

\begin{lemma}\label{atlas:crossing}
Suppose \(\overline P=\overline Q=\overline R=\overline S=0\) and
\(PS+QR=\tau c_lc_mc_tc_z\).  The following list is exhaustive.
\begin{enumerate}
\item If \(PS=0\), then for some \(a\ne0\),
\(Q=a c_lc_z\), \(R=(\tau/a)c_mc_t\), and \(P=0\) or \(S=0\).
The unused partner is any zero-double-projection tensor.
\item If \(QR=0\), then for some \(a\ne0\),
\(P=a c_lc_t\), \(S=(\tau/a)c_mc_z\), and \(Q=0\) or \(R=0\).
The unused partner is any zero-double-projection tensor.
\item If both products are nonzero, all four edges have whole rank one.
Choose one exceptional site and one row of the following table:
\[
\begin{array}{c|cccc|c}
 &P&S&Q&R&\text{condition}\\ \hline
 l&u c_t&s c_mc_z&v c_z&r c_mc_t&su+rv=\tau c_l\\
 m&p c_lc_t&u c_z&q c_lc_z&v c_t&pu+qv=\tau c_m\\
 t&c_lu&s c_mc_z&q c_lc_z&c_mv&su+qv=\tau c_t\\
 z&p c_lc_t&c_mu&c_lv&r c_mc_t&pu+rv=\tau c_z
\end{array}
\]
All named scalar parameters and both vectors \(u,v\) are nonzero.
In each row the displayed equation explicitly determines \(v\) from \(u\)
and the two scalars; the only additional restriction is \(v\ne0\).
The choices of exceptional site may overlap.
\end{enumerate}
Every listed choice is a solution.
\end{lemma}

\begin{proof}
If one product vanishes, the other equals the nonzero pure tensor.
Factorization across its own edge bipartition gives the first two cases;
a disjoint tensor product vanishes exactly when a whole factor vanishes.

Suppose both products are nonzero.  Every edge has rank at most two by
\eqref{atlas:zero-kernel}.  If \(P\) had rank two, flatten
\(PS=\tau c_lc_mc_tc_z-QR\) across \((l,z)\mid(m,t)\).
Its rank is \(\rk(P)\rk(S)\), whereas the right side has rank at most two.
Thus \(S=s_ms_z\) has rank one and the flattening has rank two.
Its left and right spaces are respectively
\[
 \im_l(P)\otimes\C s_z=\Span(c_lc_z,Q),\qquad
 \C s_m\otimes\im_t(P)=\Span(c_mc_t,R).
\]
It follows that \(S=s c_mc_z\), \(Q=q c_z\), and \(R=c_mv\), with
\(s\ne0\).  Hence \(sP=\tau c_lc_t-qv\).
Rank two of \(P\) forces \(q\notin\C c_l\) and \(v\notin\C c_t\),
contradicting \(\overline P=0\).  The same argument, with the other crossing
cut, applies to each of \(Q,R,S\).  Thus all four have rank one.

Now \(PS\) and \(QR\) are nonzero simple four-tensors whose sum is nonzero
and simple.  Their factor lines can differ at at most one site: otherwise
flattening at one differing site gives two independent factors on each side
and therefore rank two.  At every other site the common line must be the
target line \(\C c_i\).  Absorbing the scalar factors gives exactly the four
rows.  Substitution proves their converses, including the case in which all
four factor lines agree.
\end{proof}

\begin{proposition}\label{atlas:single}
Fix \(r\in\{1,2\}\).  Suppose only \(\overline L\) is nonzero and its rank
is \(r\).  Then \(K=\kappa c_tc_z\), and all solutions are the following.
\begin{enumerate}
\item If \(\kappa=0\), take \(L\) free with projected rank \(r\), and take
\(P,Q,R,S\) from Lemma~\ref{atlas:crossing}.
\item If \(\kappa\ne0\) and both crossing products are nonzero, choose
nonzero \(p,q\in V_l\) and \(r_0,s\in V_m\), and put
\begin{equation}\label{atlas:single-live}
 \begin{gathered}
 P=p c_t,\quad Q=q c_z,\quad R=r_0c_t,\quad S=s c_z,\\
 L=\kappa^{-1}(\tau c_lc_m-ps-qr_0).
 \end{gathered}
\end{equation}
Require the matrix \(M=\overline p\,\overline s+
\overline q\,\overline{r_0}\) to have rank \(r\).
Explicitly, in any quotient bases, this means
\[
\begin{cases}
 M\ne0,\quad
 \det[\overline p,\overline q]\det[\overline s,\overline{r_0}]=0,
 &r=1,\\
 \det[\overline p,\overline q]\det[\overline s,\overline{r_0}]\ne0,
 &r=2.
\end{cases}
\]
\item If \(\kappa\ne0\) and \(QR=0\), necessarily \(r=1\).
Choose \(p\in V_l,s\in V_m\) with nonzero quotients, and set
\[
 P=p c_t,\quad S=s c_z,\quad
 L=\kappa^{-1}(\tau c_lc_m-ps).
\]
Choose any zero-double-projection \(Q,R\) with \(Q=0\) or \(R=0\).
\item If \(\kappa\ne0\) and \(PS=0\), necessarily \(r=1\).
Choose \(q\in V_l,r_0\in V_m\) with nonzero quotients, and set
\[
 Q=q c_z,\quad R=r_0c_t,\quad
 L=\kappa^{-1}(\tau c_lc_m-qr_0).
\]
Choose any zero-double-projection \(P,S\) with \(P=0\) or \(S=0\).
\end{enumerate}
The branches exhaust all zero and nonzero crossing products.
\end{proposition}

\begin{proof}
Project \eqref{atlas:equation} at \(l,m,t\), and then at \(l,m,z\).
The crossing terms vanish in each case, so the nonzero \(\overline L\)
forces both partial quotients of \(K\) to vanish.  Hence \(K=\kappa c_tc_z\).
Write the other edges uniquely as
\[
 P=p c_t+c_la,\quad Q=q c_z+c_lb,\quad
 R=r_0c_t+c_md,\quad S=s c_z+c_me,
\]
where \(a,d\in H_t\), \(b,e\in H_z\), and the other vectors are whole.
Separating the four \((t,z)\) sectors gives the equivalent full system
\begin{equation}\label{atlas:single-sectors}
\begin{aligned}
 \kappa L+ps+qr_0&=\tau c_lc_m,\\
 c_lsa+qc_md&=0,&
 pc_me+c_lr_0b&=0,& ae+db&=0.
\end{aligned}
\end{equation}
When \(\kappa=0\), \(L\) drops out and Lemma~\ref{atlas:crossing} applies.

Let \(\kappa\ne0\).  The double projection of the first equation shows
that at least one of \(\overline p\,\overline s\) and
\(\overline q\,\overline{r_0}\) is nonzero.  If the first is nonzero,
projecting the second equation at \(m\) gives \(a=0\), and projecting the
third at \(l\) gives \(e=0\).  The remaining conditions are
\(qd=0\), \(r_0b=0\), and \(db=0\).  If \(Q,R\ne0\), these force
\(b=d=0\): for example \(d\ne0\) would imply \(q=b=0\), hence \(Q=0\).
If \(QR=0\), the unused partner is unrestricted within
\eqref{atlas:zero-kernel}.  Interchanging the two crossing products proves
the corresponding assertion when the second projected summand is nonzero.
Both whole products cannot vanish, since that would make \(\overline L=0\).
This proves all the stated forms.  The determinant formula follows by
writing \(M=[\overline p,\overline q][\overline s,\overline{r_0}]^T\).
A single nonzero projected outer product has rank one, explaining the last
two branches.  Finally the displayed formulas solve the entire system
\eqref{atlas:single-sectors}, which proves every converse.
\end{proof}

\subsection{Two adjacent projected edges}

Here \(\overline L,\overline P\ne0\) and
\(\overline Q=\overline R=\overline S=\overline K=0\).

\begin{lemma}\label{atlas:wedge-reduction}
Every solution in this support has
\begin{equation}\label{atlas:wedge-reduced}
 Q=q c_z,\quad S=s c_z,\quad K=k c_z,\qquad
 Lk+Ps+qR=\tau c_lc_mc_t,
\end{equation}
where \(q\ne0\) and \(R\ne0\).  Either both \(\overline s,\overline k\)
vanish or both are nonzero.
If both vanish, write \(k=\kappa c_t\), \(s=\sigma c_m\), and use the unique
decompositions
\[
 \begin{split}
 R&=r_Hc_t+c_mt_H+\rho c_mc_t,\\
 L&=L_H+xc_m,\qquad P=P_H+yc_t,
 \end{split}
\]
with \(r_H\in H_m\), \(t_H\in H_t\),
\(L_H\in V_l\otimes H_m\), \(P_H\in V_l\otimes H_t\).
Then the reduced equation is equivalent to
\begin{equation}\label{atlas:wedge-scalars}
 \kappa L_H+qr_H=0,\qquad
 \sigma P_H+qt_H=0,\qquad
 \kappa x+\sigma y+\rho q=\tau c_l.
\end{equation}
\end{lemma}

\begin{proof}
Write
\(Q=q c_z+c_lu\), \(S=s c_z+c_md\), \(K=k c_z+c_tf\), with
\(u,d,f\in H_z\).  Projection at \(l,m,z\) gives
\(\overline Lc_tf=0\), so \(f=0\); projection at \(l,t,z\) similarly gives
\(d=0\).
We claim that
\begin{equation}\label{atlas:two-term-obstruction}
 Lk+Ps=\tau c_lc_mc_t
\end{equation}
is impossible when both double projections are nonzero.
Its full quotient shows that if one of \(\overline k,\overline s\) is zero,
then both are zero.  Writing \(k=\kappa c_t,s=\sigma c_m\), projection at
\(l,m\) and \(l,t\) forces \(\kappa=\sigma=0\), a contradiction.
If both quotients are nonzero, quotient the whole spaces at \(m,t\) by
\(\C s,\C k\).  The left side vanishes and the pure target survives,
again a contradiction.
Thus \(R\ne0\).  The \(H_z\) component of the original equation is
\(c_luR=0\), giving \(u=0\).  Cancel \(c_z\) to obtain
\eqref{atlas:wedge-reduced}; the same obstruction gives \(q\ne0\).
Full projection of that equation is
\(\overline L\,\overline k+\overline P\,\overline s=0\), proving the
zero/nonzero dichotomy.  Finally the \(H_mc_t,c_mH_t,c_mc_t\) sectors are
exactly \eqref{atlas:wedge-scalars}.
\end{proof}

\begin{proposition}\label{atlas:wedge11}
For \(\rk\overline L=\rk\overline P=1\), the following five families are
disjoint and exhaustive.

\paragraph{I.} Choose \(\eta\ne0\), and take \(L,P\) free with projected
rank one.  Set
\[
 Q=\eta c_lc_z,\qquad R=(\tau/\eta)c_mc_t,\qquad S=K=0.
\]

\paragraph{II.} Choose \(\kappa\ne0\), \(q\in V_l,r\in V_m\) with
\(\overline q,\overline r\ne0\), and take \(P\) free with projected rank one.
Set
\[
 L=(\tau c_lc_m-qr)/\kappa,\quad Q=q c_z,\quad R=r c_t,
 \quad S=0,\quad K=\kappa c_tc_z.
\]

\paragraph{III.} Choose \(\sigma\ne0\), \(q\in V_l,r\in V_t\) with
\(\overline q,\overline r\ne0\), and take \(L\) free with projected rank one.
Set
\[
 P=(\tau c_lc_t-qr)/\sigma,\quad Q=q c_z,\quad R=c_mr,
 \quad S=\sigma c_mc_z,\quad K=0.
\]

\paragraph{IV.} Choose \(\kappa,\sigma\ne0\), \(q\in V_l\) with
\(\overline q\ne0\), nonzero \(A\in H_m,B\in H_t\), arbitrary
\(x\in V_l\) and \(\rho\in\C\).  Put
\[
 y=(\tau c_l-\kappa x-\rho q)/\sigma
\]
and set
\begin{equation}\label{atlas:wedge11-IV}
 \begin{gathered}
 L=qA+xc_m,\qquad P=qB+yc_t,\qquad Q=q c_z,\\
 R=-\kappa Ac_t-\sigma c_mB+\rho c_mc_t,\qquad
 S=\sigma c_mc_z,\qquad K=\kappa c_tc_z.
 \end{gathered}
\end{equation}

\paragraph{V.} Choose \(\eta\ne0\), vectors \(v\in V_l,s\in V_m,k\in V_t\)
with all three quotients nonzero, and arbitrary \(\alpha,\beta\in\C\).
Set
\begin{equation}\label{atlas:wedge11-V}
 \begin{gathered}
 L=vs+\alpha c_lc_m,\qquad P=-vk+\beta c_lc_t,\qquad Q=\eta c_lc_z,\\
 R=(\tau c_mc_t-\alpha c_mk-\beta s c_t)/\eta,\qquad
 S=s c_z,\qquad K=k c_z.
 \end{gathered}
\end{equation}
In IV and V all six whole tensors are nonzero and the two projected
\(l\)-endpoint lines coincide.  Independent \(l\)-endpoint lines occur
only in I--III, with the indicated condition imposed on the free tensors.
\end{proposition}

\begin{proof}
Use Lemma~\ref{atlas:wedge-reduction}.  If \(\overline k=\overline s=0\),
consider \eqref{atlas:wedge-scalars}.
For \(\kappa=\sigma=0\), its first two equations give \(r_H=t_H=0\), and
\(\rho q=\tau c_l\) gives I.
For \(\kappa\ne0,\sigma=0\), we have \(t_H=0\) and
\(\kappa L=\tau c_lc_m-q(r_H+\rho c_m)\).  The nonzero projected rank
forces \(\overline q,r_H\ne0\), giving II.  The opposite scalar choice
gives III.  If both scalars are nonzero, set
\(A=-r_H/\kappa\), \(B=-t_H/\sigma\).  Then the first two equations give
\(L_H=qA,P_H=qB\), with \(\overline q,A,B\ne0\), and the last equation
determines \(y\).  This is IV.

Suppose now \(\overline s,\overline k\ne0\).  Quotient
\eqref{atlas:wedge-reduced} at \(m,t\) by the actual lines \(\C s,\C k\).
The target survives, so the remaining equality \(qR'=\tau c_l[c_m][c_t]\)
forces \(q=\eta c_l\) with \(\eta\ne0\).
Projection at \(l\) gives
\(L_lk+P_ls=0\).  Quotienting at \(m\) by \(\C s\) and then canceling
the nonzero vectors shows
\(L_l=us,P_l=-uk\), for a nonzero \(u\in H_l\).
Consequently \(L=us+c_lA_m\) and \(P=-uk+c_lB_t\).
After substitution, the remaining equation is
\[
 A_mk+sB_t+\eta R=\tau c_mc_t.
\]
Its double quotient is
\(\overline{A_m}\,\overline k+\overline s\,\overline{B_t}=0\).
Because \(\overline s,\overline k\ne0\), contraction by a functional
nonzero on either factor gives
\[
 A_m=\gamma s+\alpha c_m,\qquad
 B_t=-\gamma k+\beta c_t
\]
for scalars \(\gamma,\alpha,\beta\).  Set \(v=u+\gamma c_l\).
This gives exactly V.

For completeness the converses are whole-tensor identities.  In I only
\(qR\) contributes.  In II and III its non-target part cancels the corresponding
part of \(Lk\) or \(Ps\).  In IV the sum is
\[
 (\kappa x+\sigma y+\rho q)c_mc_t=\tau c_lc_mc_t,
\]
since the terms in \(qAc_t\) and \(qc_mB\) cancel separately.
In V the terms \(vsk\), \(\alpha c_lc_mk\), and \(\beta c_ls c_t\)
each cancel, leaving the same target.  Multiplying by \(c_z\) proves
\eqref{atlas:equation}.  The prescribed projected ranks follow from the
displayed nonzero outer factors.

I--III are distinguished by which of \(S,K\) vanish.  In IV both are
nonzero with zero leaf quotients; in V both leaf quotients are nonzero.
Thus the families are disjoint.  IV and V visibly have common projected
\(l\)-endpoint lines.  Conversely, if the lines of \(\overline L\) and
\(\overline P\) are independent, the equation
\(\overline L\,\overline k+\overline P\,\overline s=0\) forces both
leaf quotients zero, and IV still has coincident lines.  This gives the
last assertion.
\end{proof}

\begin{remark}\label{atlas:wedge-parameters}
With the fixed complements, the parameters in each wedge family are
recoverable from the whole tuple.  First recover \(q,s,k\) by contraction
against the covector taking \(c_z\) to one and vanishing on \(H_z\).
In IV, the unique decomposition of \(R\) gives \(A,B,\rho\), and
contraction of \(L\) at \(c_m\) gives \(x\).  In V, contraction of \(L\)
against any covector vanishing on \(c_m\) and taking \(s\) to one recovers
\(v\); the residual pure tensors recover \(\alpha,\beta\).
The simpler recoveries in I--III follow directly from their formulas.
\end{remark}

\begin{proposition}\label{atlas:wedge21}
For \(\rk\overline L=2\), \(\rk\overline P=1\), and the other four
projections zero, \(K=0\) and \(L\) is free with projected rank two.
The remaining tensors have precisely one of the following forms:
\begin{enumerate}
\item \(S=0\), \(Q=a c_lc_z\), \(R=(\tau/a)c_mc_t\), with \(a\ne0\),
and \(P\) free with projected rank one;
\item for \(\sigma\ne0\), \(u\in V_l,v\in V_t\) with
\(\overline u,\overline v\ne0\),
\[
 S=\sigma c_mc_z,\quad Q=u c_z,\quad R=c_mv,\quad
 P=(\tau c_lc_t-uv)/\sigma.
\]
\end{enumerate}
\end{proposition}

\begin{proof}
In \eqref{atlas:wedge-reduced}, write \(\overline P=xy\) with
\(x,y\ne0\).  Full projection gives
\(\overline L\,\overline k+x\overline s\,y=0\).
A nonzero functional at \(l\) annihilating \(x\) cannot annihilate the
rank-two tensor \(\overline L\); hence \(\overline k=0\), and then
\(\overline s=0\).  In \eqref{atlas:wedge-scalars}, projection of the
first equation at \(l\) is a rank-two tensor multiplied by \(\kappa\)
plus one outer product.  Thus \(\kappa=0\), and then \(r_H=0\) because
\(q\ne0\).  If \(\sigma=0\), the remaining equations give the first
branch.  Otherwise write \(u=q\), \(v=t_H+\rho c_t\); they give the second
branch and its nonzero quotient conditions.  Its whole \(P\) has rank two,
since \((c_l,u)\) and \((c_t,v)\) are independent pairs.  Direct substitution
proves both converses, with no further restriction on \(L\).
\end{proof}

\begin{proposition}\label{atlas:wedge22}
For \(\rk\overline L=\rk\overline P=2\), and the other four projections
zero, all solutions are
\[
 S=K=0,\qquad Q=a c_lc_z,\qquad R=(\tau/a)c_mc_t,\qquad a\ne0,
\]
with \(L,P\) free with projected rank two.
\end{proposition}

\begin{proof}
Full projection of \eqref{atlas:wedge-reduced} gives
\(\overline L\,\overline k+\overline P\,\overline s=0\).
If \(\overline k\ne0\), contract at \(t\) with a nonzero functional
annihilating it.  Rank two of \(\overline P\) forces \(\overline s=0\),
and the original equality then forces \(\overline k=0\), a contradiction.
Thus both leaf quotients vanish.  Projecting the first two equations of
\eqref{atlas:wedge-scalars} at \(l\), rank two versus rank at most one
gives \(\kappa=\sigma=0\).  The remaining equations give the stated pure
\(Q,R\).  Their product is the target, proving the converse for arbitrary
whole lifts of \(L,P\).
\end{proof}

\subsection{The two triangle types}

The triangle is on \(l,m,t\), with \(z\) its isolated projected vertex.
The unique pure spoke from Lemma~\ref{support:triangle} distinguishes one
of these three vertices.  In the next statement call that vertex \(t\).
Choosing each of the three possible vertices in turn covers all labelled
triangle solutions.

\begin{proposition}\label{atlas:triangle}
Choose \(\kappa\ne0\), \(q,v\in V_l\), \(s,w\in V_m\), and
\(h\in H_t\), with \(\overline q,\overline s,h\ne0\), and set
\begin{equation}\label{atlas:triangle-form}
 \begin{gathered}
 Q=q c_z,\qquad S=s c_z,\qquad K=\kappa c_tc_z,\\
 P=qh+vc_t,\qquad R=-sh+wc_t,\qquad
 L=(\tau c_lc_m-vs-qw)/\kappa.
 \end{gathered}
\end{equation}
Extend \(\overline q\) to a quotient basis \((\overline q,e)\) at \(l\),
and \(\overline s\) to a quotient basis \((\overline s,f)\) at \(m\).
Write
\[
 \overline v=v_1\overline q+v_2e,\qquad
 \overline w=w_1\overline s+w_2f,\qquad d=v_1+w_1.
\]
The exact additional restrictions for the two triangle types are
\begin{equation}\label{atlas:triangle-test}
 \begin{array}{ll}
 (1,1,1):&v_2w_2=0,\quad(d,v_2,w_2)\ne(0,0,0),\\
 (2,1,1):&v_2w_2\ne0.
 \end{array}
\end{equation}
In the second row \(L\) is the rank-two projected edge.
These formulas exhaust both types and give their full converses.
For the fixed complement \(H_t\), their parameters are unique once the
pure-spoke vertex and the ordering of the other two vertices are fixed.
\end{proposition}

\begin{proof}
Apply Lemma~\ref{support:triangle}, with its pure-spoke vertex taken as \(t\).
It gives the whole forms \eqref{atlas:triangle-form} with a lift of a
nonzero quotient vector in the \(t\) slot.  Changing that lift by
\(h\mapsto h+\lambda c_t\) changes the remaining vectors by
\[
 v\mapsto v-\lambda q,\qquad w\mapsto w+\lambda s.
\]
These simultaneous changes leave every tensor unchanged, since
\(vs+qw\) is unchanged.  They allow and uniquely determine the choice
\(h\in H_t\).

The two incident projections are
\(\overline P=\overline q\,h\) and
\(\overline R=-\overline s\,h\), both nonzero of rank one.  The unscaled
numerator of the opposite double projection has matrix
\[
 \overline v\,\overline s+\overline q\,\overline w
 \quad\longleftrightarrow\quad
 \begin{pmatrix}d&w_2\\v_2&0\end{pmatrix}.
\]
Its determinant is \(-v_2w_2\).  A nonzero matrix of this size has rank one
exactly when its determinant vanishes, and rank two exactly when it does
not.  This proves \eqref{atlas:triangle-test}, including the branches
\(v_2=0\) and \(w_2=0\) and their intersection; the excluded zero triple
is exactly a vanished opposite projection.

Conversely the three products from \eqref{atlas:triangle-form} are
\[
\begin{aligned}
 LK&=\tau c_lc_mc_tc_z-vsc_tc_z-qwc_tc_z,\\
 PS&=qshc_z+vsc_tc_z,\\
 QR&=-qshc_z+qwc_tc_z.
\end{aligned}
\]
Their sum is the required whole target, with all scalar components retained.
The rank test gives the desired support exactly.
Finally \(Q,S,K\) recover \(q,s,\kappa\) by contraction at \(c_z\).
The partial \(t\)-projection of \(P\) recovers the quotient of \(h\),
and the fixed complement recovers \(h\) itself.  The residual tensors in
\(P,R\) then recover \(v,w\).  This proves uniqueness.
\end{proof}

\begin{remark}\label{atlas:triangle-ranks}
All six tensors and all three products are nonzero in either triangle type.
For type \((2,1,1)\), the pairs \((\overline q,\overline v)\) and
\((\overline s,\overline w)\) are independent.  In the bases
\((c_l,q,v)\), \((c_m,s,w)\), the matrix of \(L\) is
\[
 \kappa^{-1}
 \begin{pmatrix}\tau&0&0\\0&0&-1\\0&-1&0\end{pmatrix}.
\]
Thus its whole rank is three, while \(P,R\) have whole rank two and the
spokes have rank one.  For type \((1,1,1)\), either
\(\Span(c_l,q,v)\) or \(\Span(c_m,s,w)\) has dimension at most two,
so all three triangle tensors have rank at most two.  Vanishing individual
coordinates or further whole-rank drops require no additional chart.
\end{remark}

\subsection{The four-cycle}

\begin{proposition}\label{atlas:cycle}
Suppose \(\overline L=\overline K=0\), and
\(\overline P,\overline Q,\overline R,\overline S\) are nonzero of rank one.
All solutions are exactly
\begin{equation}\label{atlas:cycle-form}
 \begin{gathered}
 L=\gamma c_lc_m,\qquad K=(\tau/\gamma)c_tc_z,\qquad\gamma\ne0,\\
 P=u_lu_t,\qquad Q=u_lu_z,\qquad
 R=u_mu_t,\qquad S=-u_mu_z,
 \end{gathered}
\end{equation}
where \(u_i\in V_i\) and \(\overline{u_i}\ne0\) for all four sites.
All six whole tensors have rank one, and
\(PS+QR=0\), \(LK=\tau c_lc_mc_tc_z\).
\end{proposition}

\begin{proof}
Full projection gives
\(\overline P\,\overline S=-\overline Q\,\overline R\ne0\).
Equality of nonzero simple four-tensors gives equal endpoint lines at
every site.  The scalar factors can be chosen consistently: first factor
\(\overline P=a_la_t\), then \(\overline Q=a_la_z\), then
\(\overline R=a_ma_t\); the equality forces
\(\overline S=-a_ma_z\).  Here each \(a_i\) is a nonzero vector in \(H_i\).
Retain all whole lifts by writing
\[
\begin{aligned}
 P&=a_la_t+p_lc_t+c_lp_t+p_0c_lc_t,\\
 Q&=a_la_z+q_lc_z+c_lq_z+q_0c_lc_z,\\
 R&=a_ma_t+r_mc_t+c_mr_t+r_0c_mc_t,\\
 S&=-a_ma_z+s_mc_z+c_ms_z+s_0c_mc_z,
\end{aligned}
\]
with the subscripted vector at site \(i\) in \(H_i\).
Since \(\overline L=\overline K=0\), their product contributes no
three-\(H\) component.  The four three-\(H\) components of the full equation
therefore give
\[
 a_ls_m+q_la_m=0,\quad -p_la_m+a_lr_m=0,\quad
 a_ts_z+r_ta_z=0,\quad -p_ta_z+a_tq_z=0.
\]
For nonzero \(a,b\), an equality \(ax+yb=0\) forces
\(x=\lambda b,y=-\lambda a\), including \(\lambda=0\): quotienting by
either fixed line and then canceling the other factor proves this.
Consequently there are scalars \(\alpha,\beta,\delta,\epsilon\) with
\[
\begin{gathered}
 p_l=\alpha a_l,\quad r_m=\alpha a_m,\quad
 q_l=-\beta a_l,\quad s_m=\beta a_m,\\
 p_t=\delta a_t,\quad q_z=\delta a_z,\quad
 r_t=\epsilon a_t,\quad s_z=-\epsilon a_z.
\end{gathered}
\]
Define whole vectors
\[
 x_l=a_l+\delta c_l,\quad x_m=a_m+\epsilon c_m,\quad
 x_t=a_t+\alpha c_t,\quad x_z=a_z-\beta c_z.
\]
Expansion gives
\begin{equation}\label{atlas:cycle-residual}
 \begin{gathered}
 P=x_lx_t+A c_lc_t,\qquad Q=x_lx_z+B c_lc_z,\\
 R=x_mx_t+C c_mc_t,\qquad S=-x_mx_z+D c_mc_z,
 \end{gathered}
\end{equation}
where
\(A=p_0-\delta\alpha\), \(B=q_0+\delta\beta\),
\(C=r_0-\epsilon\alpha\), \(D=s_0-\epsilon\beta\).
All \(\overline{x_i}\) are nonzero.

Choose complements to \(\C c_i\) containing \(x_i\); this is merely a
choice of basis for the following coefficient comparison, and does not
change the tensors or \(\tau\).  Write the still unrestricted missing edges
in these complements as
\[
 L=\xi_lc_m+c_l\eta_m+\ell c_lc_m,\qquad
 K=\rho_tc_z+c_t\sigma_z+\kappa c_tc_z.
\]
The four cross-pair components, the four single-site components, and the
constant component of \eqref{atlas:equation} are respectively
\begin{equation}\label{atlas:cycle-coefficients}
\begin{gathered}
 \xi_l\rho_t=-D x_lx_t,\qquad
 \xi_l\sigma_z=-C x_lx_z,\\
 \eta_m\rho_t=-B x_mx_t,\qquad
 \eta_m\sigma_z=A x_mx_z,\\
 \kappa\xi_l=\kappa\eta_m=\ell\rho_t=\ell\sigma_z=0,\qquad
 \ell\kappa+AD+BC=\tau.
\end{gathered}
\end{equation}
These are all remaining coefficients: the other two pair components are
identically zero, the three-site components have just been eliminated,
and the four-site components cancel in \eqref{atlas:cycle-residual}.

Multiply the first and fourth pair equations in
\eqref{atlas:cycle-coefficients}, and then the second and third, restoring
both to the same site order.  Their left sides are the identical tensor
\(\xi_l\eta_m\rho_t\sigma_z\).  Since \(x_lx_mx_tx_z\ne0\), their right
sides give \(-DA=CB\), so \(AD+BC=0\).
This argument divides by none of the unknown vectors or residual scalars.
The final equation now gives \(\ell\kappa=\tau\ne0\).
The single-site equations force \(\xi_l=\eta_m=\rho_t=\sigma_z=0\),
and the pair equations force \(A=B=C=D=0\).
Taking \(\gamma=\ell\) and \(u_i=x_i\) yields
\eqref{atlas:cycle-form} in the original spaces.

Conversely let the parameters in \eqref{atlas:cycle-form} be given, and put
\(U=u_lu_mu_tu_z\).  Then \(LK=\tau c_lc_mc_tc_z\), \(PS=-U\), and
\(QR=U\).  The double projections have exactly the prescribed ranks,
and every whole factor is nonzero of rank one.  This proves the full
converse, including arbitrary \(c\)-components in the four vectors.
\end{proof}

The propositions describe whole tensors, rather than only quotient data.
Every displayed parametrization therefore includes all permitted cancellations
and whole-rank drops.  The list is finite at the level of support and rank
types; Proposition~\ref{atlas:empty} uses geometric parameters, and no
classification into finitely many orbits is asserted.
